% This must be in the first 5 lines to tell arXiv to use pdfLaTeX, which is strongly recommended.
% \pdfoutput=1
% In particular, the hyperref package requires pdfLaTeX in order to break URLs across lines.

\documentclass[11pt]{article}

% Remove the "review" option to generate the final version.
\usepackage{ACL2023}

% Standard package includes
\usepackage{times}
\usepackage{latexsym}

% For proper rendering and hyphenation of words containing Latin characters (including in bib files)
\usepackage[T1]{fontenc}
% For Vietnamese characters
% \usepackage[T5]{fontenc}
% See https://www.latex-project.org/help/documentation/encguide.pdf for other character sets

% This assumes your files are encoded as UTF8
\usepackage[utf8]{inputenc}

% This is not strictly necessary, and may be commented out.
% However, it will improve the layout of the manuscript,
% and will typically save some space.
\usepackage{microtype}

% This is also not strictly necessary, and may be commented out.
% However, it will improve the aesthetics of text in
% the typewriter font.
\usepackage{inconsolata}


% If the title and author information does not fit in the area allocated, uncomment the following
%
%\setlength\titlebox{<dim>}
%
% and set <dim> to something 5cm or larger.

\title{The Illusion of Objectivity in RAG Evaluation: A Survey on the Reliability and Biases of LLM Judges}

% Author information can be set in various styles:
% For several authors from the same institution:
% \author{Author 1 \and ... \and Author n \\
%         Address line \\ ... \\ Address line}
% if the names do not fit well on one line use
%         Author 1 \\ {\bf Author 2} \\ ... \\ {\bf Author n} \\
% For authors from different institutions:
% \author{Author 1 \\ Address line \\  ... \\ Address line
%         \And  ... \And
%         Author n \\ Address line \\ ... \\ Address line}
% To start a seperate ``row'' of authors use \AND, as in
% \author{Author 1 \\ Address line \\  ... \\ Address line
%         \AND
%         Author 2 \\ Address line \\ ... \\ Address line \And
%         Author 3 \\ Address line \\ ... \\ Address line}

\author{Lynn Zhou \\
  \texttt{zzzzzzzl@tamu.edu} \And
  Deepmoy Hazra\\
  \texttt{deepmoy3@tamu.edu} \And
  Huang, Wei-Chen \\
  \texttt{wilson0623@tamu.edu} \AND
  Texas A\&M University \\
  College Station, TX 77843
}

\begin{document}
\maketitle
\begin{abstract}
Retrieval-Augmented Generation (RAG) systems increasingly rely on Large Language Models (LLMs) as automated evaluators due to the unscalable nature of human annotation. However, the reliability of these judges under the unique stressors of RAG environments remains under-explored. This proposal outlines a comprehensive survey investigating the micro-level cognitive vulnerabilities of LLM evaluators. We aim to bridge the gap between macro-level evaluation metrics and internal model mechanisms by dissecting the "Truthfulness vs. Faithfulness" paradox, analyzing how contextual noise exacerbates inherent biases, and tracking the evolutionary trajectory of mitigation strategies. Ultimately, this survey will map the critical paradigm shift from static multi-agent debates to proactive, tool-augmented "Agent-as-a-Judge" frameworks.
\end{abstract}

\section{Introduction and Scope}
Evaluating complex generative systems is notoriously resource-intensive, making the industry heavily dependent on LLMs to serve as scalable, automated judges. While this approach is highly efficient, it introduces a critical vulnerability: if the evaluator itself is compromised by contextual distractors or positional biases, the foundational reliability of the entire system degrades. The scope of our proposed survey focuses strictly on these micro-level mechanistic conflicts induced by RAG-specific environments. We aim to investigate how extreme context lengths and retrieved noise act as specific catalysts for inherent biases, and how models navigate the fundamental tension between factual correctness and contextual adherence. This analysis is vital for researchers and machine learning engineers seeking to design robust, bias-resistant assessment architectures.


\section{Related Work and Identified Gaps}
In fact, recent literature has begun to address the concept of automated evaluation. Foundational surveys, such as the works by \citet{gu2025surveyllmasajudge} and \citet{li-etal-2025-generation}, outline general evaluator vulnerabilities, including positional and verbosity biases. Furthermore, recent highly relevant work like \citet{Brehme_2025} macroscopically catalogs various evaluation methodologies and datasets specific to RAG frameworks. However, these existing overviews tend to treat RAG evaluation either as a macro-level systems engineering issue or as a standard assessment task. 

Our proposed survey aims to address three distinct limitations in the current literature. First, rather than broadly observing hallucination phenomena, we will dissect the internal epistemological conflict, specifically analyzing how judges prioritize parametric knowledge over contextual evidence. Second, we plan to isolate the specific intersection where RAG-induced noise exacerbates the vulnerability of the foundational evaluator. Finally, we will critically map the limitations of current static consensus solutions, using these observed failures to bridge the transition toward active, multi-step agentic verification methods as recently pioneered by \citet{you2026agentasajudge}.

\section{Proposed Survey Structure}
The survey will be organized following a mechanistic and evolutionary principle, strategically designed to build a cohesive narrative from problem identification to paradigm shift. We will begin by establishing the environmental stressors unique to RAG systems, then open the "black box" of the LLM evaluator to understand the internal cognitive conflicts these stressors provoke, and finally trace the trajectory of solutions developed by the research community. To achieve this logical flow, the survey is divided into three interconnected modules:

\begin{itemize}
\item \textbf{Phenomena and Catalysts:} This section will analyze how the unique stressors of retrieval frameworks, specifically extreme context lengths and the introduction of retrieved distractors, trigger and amplify evaluator bias.
\item \textbf{Core Conflict Mechanisms:} Building upon the identified environmental triggers, this module will provide a deep dive into the internal pathology of the judge. It focuses on the truthfulness versus faithfulness paradox and examines how latent knowledge conflicts are processed within the neural architecture of the evaluator.
\item \textbf{Mitigation and Paradigm Shifts:} Concluding the narrative arc, the final section will track the evolutionary progression of calibration strategies. We will trace the movement from static interventions, like chain-of-thought prompting and consensus juries, to the ultimate necessity of proactive, tool-augmented agentic frameworks capable of external verification.
\end{itemize}

\section{Expected Insights}
By synthesizing the latest advancements, readers will gain a granular understanding of why standard LLM judges are highly susceptible to contextual manipulation and cannot be implicitly trusted in high-noise environments. Furthermore, the survey will demonstrate that forcing agreement through multi-model debates often results in biased compromises rather than authentic error correction in complex scenarios. Ultimately, the primary insight we emphasize is the necessity of a paradigm shift within the field, moving from passive evaluators to active, investigating agents that utilize multi-step verification and external tool invocation to resolve complex knowledge conflicts.

\section{Project Timeline}
\section{Project Timeline}
To align with the presentation schedule beginning April 20, we have established a highly parallelized and strictly milestone-driven timeline:
\begin{itemize}
    \item{\textbf{Mar 5–15: Literature Search \& Framework Detailing.}} Finalize the comprehensive corpus of core papers and establish the granular structural hierarchy for the entire survey. 
    \item{\textbf{Mar 16 – Apr 5: Concurrent Drafting of Three Modules.}} Execute parallel drafting of the three distinct sections, focusing simultaneously on RAG-specific triggers, conflict mechanisms, and mitigation paradigms. 
    \item{\textbf{Apr 6–12: Integration \& Logical Refinement.}} Merge the parallel drafts into a cohesive manuscript, unifying the academic tone, ensuring seamless transitional flow across modules, and formatting citations. 
    \item{\textbf{Apr 13–19: Presentation Preparation.}} Synthesize the core findings to develop presentation slides and conduct rehearsal sessions to distill complex insights for an audience.
    \item{\textbf{Apr 20 Onwards: In-Class Presentations \& Final Polish.}} Deliver the survey findings in class, utilizing subsequent feedback to finalize the manuscript for formal submission.
\end{itemize}

% Entries for the entire Anthology, followed by custom entries
\bibliography{anthology,custom}
\bibliographystyle{acl_natbib}



\end{document}
